\section{ARCH-15-HONDA-CHECK: full check of CC-AG14 and the class-number counterexample}
\label{arch15-honda-check}

\paragraph{Nature and scope of the check.}
This is a targeted whole-proof check of the supplied CC-AG14, including
its seven numbered proof parts and its classical abelian-variety
corollary, together with the use of that corollary in
BR-15-CLASSBOUND. It is not claimed as a proof independently discovered
without reading the checkpoint. The class-number calculation was
already confirmed by the intake; it is recalled only to connect the
verified index to the stated counterexample. The same reading also
checked BR-15-SLOPEBOUND. No mathematical error was found in these
arguments. The classical conclusions use Honda--Tate and elementary
number theory only, independently of the checkpoint's additional
higher-cycle inputs. Tags: published theorem (Honda--Tate and Tate's
isogeny criterion), none (the explicit arithmetic); A, M, and all three
named outside CM preprints are unused.

\subsection{Index, dimension, and geometric simplicity}

Write
\[
 K=\mathbf Q(s),\quad s^2=-23,\quad
 \omega=(1+s)/2,\quad \eta=1+\omega,
 \quad\gamma=8\bar\eta=12-4s.
\]
The ring \(\mathcal O_K=\mathbf Z[\omega]\) has discriminant \(-23\),
and \(N(x+y\omega)=x^2+xy+6y^2\). Set
\[
 \mathfrak p=(2,\omega-1),\qquad
 \bar{\mathfrak p}=(2,\omega).
\]
The polynomial \(t^2-t+6\) has two distinct roots modulo 2. Thus
\((2)=\mathfrak p\bar{\mathfrak p}\), both primes have norm 2, and
both completions are \(\mathbf Q_2\). The residue and norm calculations
give
\[
 (\eta)=\mathfrak p^3,\qquad
 (\gamma)=\mathfrak p^3\bar{\mathfrak p}^{6},\qquad
 \gamma\bar\gamma=512,
 \quad\gamma+\bar\gamma=24.
 \tag{ARCH-15.H1}
\]
No element has norm 2: the expression
\((x+y/2)^2+23y^2/4\) excludes \(y\ne0\), and \(y=0\) would require
\(x^2=2\). Hence \([\mathfrak p]\) has order exactly 3.

The number \(\gamma\) is an integral Weil \(2^9\)-number. For the
endomorphism division algebra of its Honda--Tate abelian variety,
Tate's positive-sign convention gives
\[
 \operatorname{inv}_{\mathfrak p}(D)=\frac39[ K_{\mathfrak p}:\mathbf Q_2]
     =\frac13,
 \qquad
 \operatorname{inv}_{\bar{\mathfrak p}}(D)=\frac69
     [K_{\bar{\mathfrak p}}:\mathbf Q_2]=\frac23
     \quad\text{in }\mathbf Q/\mathbf Z.
 \tag{ARCH-15.H2}
\]
All other finite invariants vanish; there are no real places. The
global index is therefore 3, not 1. Since \([K:\mathbf Q]=2\),
\[
 2\dim A=[K:\mathbf Q]\sqrt{[D:K]}=2\cdot3=6,
 \qquad f_A(t)=(t^2-24t+512)^3.
 \tag{ARCH-15.H3}
\]
These are exactly Theorem 1(ii), equations (1)--(2), and Remark 2 of
Tate's Bourbaki expose 352, printed p.~96. That primary scan, including
the displayed formulas, was visually inspected in this run. The
cohomological-motive convention in CC-AG14 has the opposite signs in
(ARCH-15.H2); the denominators and index are unchanged.

For every \(n\ge1\), the valuations of \(\gamma^n\) at the two primes
are \(3n,6n\), so \(\gamma^n\notin\mathbf Q\) and
\(\mathbf Q(\gamma^n)=K\). Relative to \(2^{9n}\) the invariants
are still \(1/3,2/3\). The simple variety associated to \(\gamma^n\)
therefore still has dimension 3, and Tate's polynomial isogeny
criterion identifies it with \(A_{\mathbf F_{2^{9n}}}\). Thus \(A\)
is simple after every finite extension and is geometrically simple.
Every geometric endomorphism is defined over a finite extension, while
the endomorphism algebras over those fields have the same dimension
\(2\cdot3^2=18\) over \(\mathbf Q\). Their injective transition maps
are isomorphisms. The geometric endomorphism algebra is consequently
the same \(D\), with imaginary quadratic center \(K\) and degree 3.
This checks the type \(\mathrm{IV}(1,3)\) hypothesis, not just the
Frobenius polynomial.

\subsection{Why no single geometric copy descends to \(\mathbf F_8\)}

First check the stabilized exponent used in CC-AG14. If
\(\delta\in K\) represents the same germ as \(\gamma\) at exponent
\(n\), then for some \(N>0\),
\(\delta^{9N}=\gamma^{nN}\). Therefore
\[
 \operatorname{ord}_{\mathfrak p}\delta=n/3,\qquad
 \operatorname{ord}_{\bar{\mathfrak p}}\delta=2n/3.
\]
Integrality of these valuations forces \(3\mid n\), and principality
of the resulting ideal forces \(3\mid n/3\), because
\([\mathfrak p]\) has order 3. Since \(\gamma\) exists at exponent
9, the minimum is exactly \(h_\xi=9\).

Suppose \(C/\mathbf F_8\) were geometrically isogenous to \(A\).
It would be a geometrically simple threefold. Let \(\alpha\) be its
Weil label and \(L=\mathbf Q(\alpha)\). After conjugation and passage
to a common finite field,
\[
 \alpha^{9N}=\gamma^{3N}
 \tag{ARCH-15.H4}
\]
for some \(N>0\). The right side generates \(K\), hence \(K\subset L\).
Put \(c=[L:K]\). The norm of (ARCH-15.H4) shows that
\(N_{L/K}\alpha\) represents the same germ at exponent \(3c\), so
\(3\mid c\). There are no real places of \(L\). At a place \(z\mid2\),
with integer-normalized valuation and residue degree \(f_z\), the
classical local invariant is
\[
 \frac{\operatorname{ord}_z\alpha}{\operatorname{ord}_z8}
          [L_z:\mathbf Q_2]
 =\frac{\operatorname{ord}_z\alpha\,f_z}{3}
          \pmod {\mathbf Z}.
 \tag{ARCH-15.H5}
\]
Thus the index \(d_\alpha\) is 1 or 3. The dimension formula for a
threefold says \(c d_\alpha=3\), forcing \(c=3\), \([L:\mathbf Q]=6\),
and \(d_\alpha=1\).

The full root-of-unity possibility must still be eliminated.
Equation (ARCH-15.H4) shows \(\alpha^3/\gamma\in\mu(L)\). A
degree-six field containing \(\mathbf Q(\sqrt{-23})\) has only
\(\pm1\) as roots of unity. Indeed \(\varphi(m)\mid6\) restricts the
orders to \(1,2,3,4,6,7,9,14,18\). Orders divisible by 3 or 4 would
give a second quadratic subfield and a degree-four compositum inside
\(L\), impossible. Orders 7 or 14 would give
\(L=\mathbf Q(\zeta_7)\), whose unique quadratic subfield is
\(\mathbf Q(\sqrt{-7})\). The elementary period calculation
\((\zeta_7+\zeta_7^2+\zeta_7^4)^2+
(\zeta_7+\zeta_7^2+\zeta_7^4)+2=0\) verifies that field explicitly.
After replacing \(\alpha\) by \(-\alpha\), we must therefore have
\(\alpha^3=\gamma\).

Every 2-adic unit is a cube, by Hensel lifting the simple root 1
modulo 2 of \(t^3-u\). The two valuations of \(\gamma\), namely 3 and
6, are divisible by 3. Hence \(\gamma\) has a cube root
\(b_v\in K_v=\mathbf Q_2\) at both places. The minimal polynomial
\(t^3-\gamma\) of \(\alpha\) over \(K\) has the factor
\(t-b_v\) locally, giving a degree-one local place of \(L\) with
\(\operatorname{ord}_2\alpha=1\), respectively 2. Formula
(ARCH-15.H5) now gives a nonzero invariant \(1/3\), respectively
\(2/3\), so \(d_\alpha=3\), a contradiction. This checks the
exclusion for all possible labels and all root-of-unity twists.

\subsection{The two witnesses and the complete multiplicity monoid}

For \(\alpha_0^3=\gamma\), the norm constraint makes
\([\mathbf Q(\alpha_0):K]=3\). The local linear factors just used
force index 3. Thus its Honda--Tate variety over \(\mathbf F_8\) has
dimension 9. Its center polynomial is
\[
 P_0(t)=t^6-24t^3+512.
\]
For \(\alpha_1=\alpha_0\zeta_9\),
\(\alpha_1^9=\gamma^3\) and \(\alpha_1^3/\gamma=\zeta_3\).
The relative center degree is divisible by both 2 and 3, and the
annihilating polynomial \(t^6+\gamma t^3+\gamma^2\) gives upper bound
6; thus it is exactly 6. Locally, scaling by \(b_v\) changes this
polynomial to \(\Phi_9(t)=t^6+t^3+1\). Its reduction is irreducible
over \(\mathbf F_2\), since \(\operatorname{ord}_9(2)=6\). The local
degree is 6 above each place of \(K\) at 2, so the invariants are
\(6/3=2\) and \(12/3=4\), both zero. Its index is 1 and its abelian
dimension is 6. Its center polynomial is
\[
 P_1(t)=t^{12}+24t^9+64t^6+12288t^3+262144.
\]

Upon extension \(\mathbf F_8\subset\mathbf F_{2^{27}}\), Frobenius
is raised to the ninth power and both labels become
\(\beta=\gamma^3\). The target threefold remains simple with index 3.
The dimensions, or directly the characteristic polynomials, give
multiplicities 3 and 2 respectively. Further,
\[
 P_0(t)P_1(t)
 =(t^9-\gamma^3)(t^9-\bar\gamma^3)
 =t^{18}+23040t^9+134217728.
 \tag{ARCH-15.H6}
\]
The degree arguments above prove the irreducibility of both displayed
center polynomials. Hence these are exactly the two finite source
orbits for that fixed finite target, not merely the results of a
bounded search. For the geometric target, multiplicity one was
excluded without a search bound and every \(r\ge2\) is a nonnegative
combination of 2 and 3. Both descent monoids are
\(\langle2,3\rangle=\{0\}\cup\{r\ge2\}\), exactly as CC-AG14 states.

\subsection{Connection with BR-15-CLASSBOUND and the source error}

The already confirmed class-number calculation gives \(h_K=3\): the
Minkowski bound \(2\sqrt{23}/\pi<4\) leaves ideals of norms 1, 2,
and 3, and \((\omega)=\bar{\mathfrak p}\mathfrak q\) identifies the
prime classes over 3 with the two nontrivial classes generated by
\([\mathfrak p]\). Thus the index calculation and the geometric
non-descent just checked apply to a type \(\mathrm{IV}(1,3)\)
threefold with center of class number 3.

Thakur's arXiv:2001.00210v1, Definition 2.3 and Proposition 2.9 on
p.~6, asks for a model up to isogeny over \(\mathbf F_{p^{h_K}}\).
Here that field is \(\mathbf F_8\), and no such geometric isogeny
exists. The inspected v1 proof chooses a generator of
\(\mathfrak p^{h_K}\) and assigns its Honda--Tate variety the original
division algebra and dimension. That generator instead has slopes 1
and 0 at the two split primes, hence index 1 and abelian dimension 1.
This identifies the failed inference while the preceding all-label
argument refutes the conclusion itself. Nothing here addresses other
statements or later versions of that preprint.

BR-15-SLOPEBOUND was also read through. Its sufficient geometric field
\(\mathbf F_{p^{td}}\), where \(t\) is the order of the split prime
class and \(j/d\) is the required slope, retains the missing slope
factor: \(\rho=\eta^j\bar\eta^{d-j}\),
\((\eta)=\mathfrak p^t\). It has the specified local invariants and
index \(d\). Equality of ideals makes the comparison with the original
Frobenius a root of unity, removable by finite extension. The exclusion
of equal slopes is necessary to retain the imaginary quadratic
geometric center. The statement is a sufficient geometric bound, not a
minimal-field or fixed-Frobenius descent assertion.

\paragraph{Exact computational support.}
The supplied script
\nolinkurl{consequences_criteria_noncyclic_check.py} was executed
successfully in this run. It checks the quadratic reductions, both
norm-polynomial expansions, (ARCH-15.H6), the complete finite list of
cyclotomic orders with \(\varphi(m)\mid6\), the quadratic period for
\(\zeta_7\), \(\operatorname{ord}_9(2)=6\), and exact irreducibility
of \(\Phi_9\) over \(\mathbf F_2\) and \(P_0,P_1\) over
\(\mathbf Q\). It does not verify the Honda--Tate theorem or replace
the unbounded field and local-invariant argument. Its observed stdout
is preserved in \texttt{cote15\_honda\_computation.txt}.

\paragraph{Primary locators.}
Tate, expose 352, delivered November 1968 and published 1971,
Theorem 1(i)--(ii) and Remark 2, printed p.~96, scan PDF p.~3;
Tate (1966), Theorem 1(c), printed p.~139, scan PDF p.~6;
Thakur, arXiv:2001.00210v1 (the inspected PDF header is dated
1 January 2020), Definition 2.1 p.~3 and Definition 2.3 and
Proposition 2.9 with proof p.~6. The new check event retains the
checkpoint's statement IDs and does not relabel their proofs as newly
discovered results.
